\documentclass[pdflatex,sn-mathphys-num]{sn-jnl}

\usepackage{enumitem}
\usepackage{graphicx}%
\usepackage{multirow}%
\usepackage{amsmath,amssymb,amsfonts}%
\usepackage{amsthm}%
\usepackage{mathrsfs}%
\usepackage[title]{appendix}%
\usepackage{xcolor}%
\usepackage{textcomp}%
\usepackage{manyfoot}%
\usepackage{booktabs}%
\usepackage{algorithm}%
\usepackage{algorithmicx}%
\usepackage{algpseudocode}%
\usepackage{listings}
\usepackage{aliascnt}
\usepackage{slashed}

\renewcommand{\sectionautorefname}{Section}

%% as per the requirement new theorem styles can be included as shown below
\theoremstyle{thmstyleone}
\newtheorem{theorem}{Theorem}

% Lemma
\newaliascnt{lemma}{theorem}
\newtheorem{lemma}[lemma]{Lemma}
\aliascntresetthe{lemma}
\providecommand*{\lemmaautorefname}{Lemma}

% Axiom
\newaliascnt{axiom}{theorem}
\newtheorem{axiom}[axiom]{Axiom}
\aliascntresetthe{axiom}
\providecommand*{\axiomautorefname}{Axiom}

% Corollary
\newaliascnt{corollary}{theorem}
\newtheorem{corollary}[corollary]{Corollary}
\aliascntresetthe{corollary}
\providecommand*{\corollaryautorefname}{Corollary}

% Proposition
\newaliascnt{proposition}{theorem}
\newtheorem{proposition}[proposition]{Proposition}
\aliascntresetthe{proposition}
\providecommand*{\propositionautorefname}{Proposition}

\theoremstyle{thmstyletwo}%
\newtheorem{example}{Example}%
\newtheorem{remark}{Remark}%

\theoremstyle{thmstylethree}%
\newtheorem{definition}{Definition}%

\raggedbottom
%%\unnumbered% uncomment this for unnumbered level heads

% custom commands
\newcommand{\SU}[1]{\mathrm{SU(#1)}}
\newcommand{\SMSubgroup}{\SU{3}\times\SU{2}\times\U{1}}
\newcommand{\SO}[1]{\mathrm{SO(#1)}}
\newcommand{\U}[1]{\mathrm{U(#1)}}
\newcommand{\E}[1]{\mathrm{E_{#1}}}
\newcommand{\M}{\mathcal{M}}
\newcommand{\Oct}{\mathbb{O}}
\newcommand{\jalg}{J_3(\mathbb{O})}
\newcommand{\SOTen}{\SO{10}}
\newcommand{\UOne}{\U{1}}
\newcommand{\ESix}{\E{6}}
\newcommand{\SOTenXUOne}{\SOTen \times \UOne}
\newcommand{\MCoset}{\frac{\ESix}{\SOTenXUOne}}
\newcommand{\MCosetInline}{\ESix/(\SOTenXUOne)}
\newcommand{\ESixBreaking}{\ESix \;\to\; \SOTenXUOne}
\newcommand{\map}[3]{#1 \ : \ \#2 \;\to\; \#3}
\newcommand{\set}[2]{\{ #1 \,:\, #2 \}}
\newcommand{\Herm}[2]{\mathrm{Herm}_{#1}(#2)}
\newcommand{\chr}[1]{\mathrm{char}(#1)}
\newcommand{\R}{\mathbb{R}}
\newcommand{\C}{\mathbb{C}}
\newcommand{\Tr}{\operatorname{Tr}}
\newcommand{\s}{\hspace{.1em}}
\newcommand{\sm}{\hspace{.065em}}
\newcommand{\Str}{\operatorname{Str}}
\newcommand{\Spin}{\operatorname{Spin}}
\newcommand{\appref}[1]{\hyperref[#1]{Appendix~\ref*{#1}}}
\newcommand{\subsecref}[1]{\hyperref[#1]{Section~\ref*{#1}}}
\newcommand{\smexp}[1]{^{\sm #1}}
\newcommand{\Hess}[1]{\operatorname{Hess}}
\newcommand{\dv}[1]{\mathrm{d}\smexp{#1}}

\renewcommand{\arraystretch}{1.25}

\begin{document}

\title[Universal Coset Field Theory]{Universal Coset Field Theory}

\author*[1]{\fnm{Brandon} \sur{Belna}}\email{bbelna@iu.edu}
\affil*[1]{\orgdiv{Department of Mathematics}, \orgname{Indiana University East}, \orgaddress{\street{2325 Chester Blvd}, \city{Richmond}, \postcode{47374}, \state{IN}, \country{USA}}}

\abstract{
TODO}

\keywords{TODO}

%%\pacs[JEL Classification]{D8, H51}

%%\pacs[MSC Classification]{35A01, 65L10, 65L12, 65L20, 65L70}

\maketitle

%======================================================================%
\section{Introduction}
%======================================================================%

%======================================================================%
\section{Axioms, the Albert Algebra, and Spontaneous Symmetry Breaking}
\label{sec:AxiomsAlbertSSB}
%======================================================================%
In this section, we introduce our two fundamental axioms, recall the definition and key properties of the Albert algebra $J_{3}(\mathbb O)$ and its $\ESix$ symmetry (with proofs in \appref{app:Albert}), exhibit a purely algebraic polynomial potential whose minima break $\ESix\to \SOTenXUOne$, and count the resulting $32$ Goldstone directions.

%%%%%%%%%%%%%%%%%%%%%%%%%%%%%%%%%%%%%%%%%%%%%%%%%%%%%%%%%%%%%%%%%%%%%%
\subsection{Axioms}\label{subsec:Axioms}
%%%%%%%%%%%%%%%%%%%%%%%%%%%%%%%%%%%%%%%%%%%%%%%%%%%%%%%%%%%%%%%%%%%%%%
\begin{axiom}[Maximally symmetric vacuum]\label{ax:AxiomI}
The universe begins in a unique state $\Phi\in\jalg$ that is fixed by every element of its structure group.
\end{axiom}
\begin{axiom}[Universal wavefunction]\label{ax:AxiomII}
Physical states $\psi\in\mathcal H$ evolve through a one-parameter family
$\{U(t)\}_{t\in\mathbb R}$ of unitary operators satisfying $U(0)=\mathrm{Id}$, $U(t+s)=U(t)U(s)$, and strong continuity.
\end{axiom}

\begin{remark}
\autoref{ax:AxiomII} provides the kinematical Hilbert space and the one-parameter
family of unitary operators.  Whenever we later write a functional integral, we treat it as a formal generator of connected Green functions; its constructive definition follows from the Osterwalder-Schrader reconstruction theorem applied to the Euclidean correlation functions \cite{GlimmJaffe:1987}.
\end{remark}


No additional axioms are introduced; all further constructions follow as definitions or theorems derived from these two axioms.

%%%%%%%%%%%%%%%%%%%%%%%%%%%%%%%%%%%%%%%%%%%%%%%%%%%%%%%%%%%%%%%%%%%%%%
\subsection{The Albert Algebra}
\label{subsec:Albert}
%%%%%%%%%%%%%%%%%%%%%%%%%%%%%%%%%%%%%%%%%%%%%%%%%%%%%%%%%%%%%%%%%%%%%%

Proofs for this subsection are provided in \appref{app:Albert}.

\begin{definition}[Jordan algebra {\cite{Jacobson:1968}}]
Let $F$ be a field with $\operatorname{char} F \neq 2$.
A \emph{Jordan algebra} is a commutative algebra $(A,\circ)$ satisfying
\[
  a\circ(b\circ a\smexp{2}) \;=\;(a\circ b)\circ a\smexp{2},
  \quad a,b\in A.
\]
\end{definition}

\begin{definition}[Albert algebra {\cite{Albert:1934,McCrimmon:2004}}]
The \emph{Albert algebra} is
\[
  \jalg=\{X\in M_{3}(\Oct) : X=X\smexp{\dagger}\},
  \quad
  X\circ Y=\frac12\,(XY+YX).
\]
\end{definition}

\begin{theorem}\label{thm:JordanDim}
The Albert algebra is a $27$-dimensional real Jordan algebra.
\end{theorem}

\begin{theorem}[Quadratic and cubic invariants]\label{thm:Invs}
$\jalg$ admits
\[
  Q(X)=\Tr(X\smexp{2}),
  \quad
  N(X)=\det_{\Oct}(X),
\]
both preserved by its structure group.
\end{theorem}

\begin{theorem}[Structure group {\cite{Springer:2000}}]\label{thm:E6}
The structure group of $\jalg$ is the simply connected exceptional Lie group $\ESix$.
\end{theorem}

\begin{theorem}[Orbit classification]\label{thm:OrbitClass}
For $X,Y\in\jalg$,
\[
  \bigl(Q(X),N(X)\bigr)=\bigl(Q(Y),N(Y)\bigr)
\]
if and only if there exists $g\in\ESix$ such that $g\cdot X = Y$.
\end{theorem}

\begin{corollary}[Existence of an $\SOTenXUOne$ stabiliser]
\label{cor:Isotropy}
There exists $X_{0}\in J_{3}(\mathbb{O})$ with
$\mathrm{Stab}_{\ESix}(X_{0})\cong \SOTen\times \UOne$.
\end{corollary}

%%%%%%%%%%%%%%%%%%%%%%%%%%%%%%%%%%%%%%%%%%%%%%%%%%%%%%%%%%%%%%%%%%%%%%
\subsection{Polynomial Symmetry-Breaking Potential}\label{subsec:PolyPot}
%%%%%%%%%%%%%%%%%%%%%%%%%%%%%%%%%%%%%%%%%%%%%%%%%%%%%%%%%%%%%%%%%%%%%%
Set
\[
  I_{2}(X)=Q(X),
  \quad
  I_{3}(X)=N(X),
  \quad
  c_{2}=I_{2}(X_{0}),
  \quad
  c_{3}=I_{3}(X_{0}).
\]
In particular, each simultaneous level set $\{I_{2}=c_{2},\,I_{3}=c_{3}\}$ is exactly one $\ESix$-orbit, so by choosing $c_{2}=I_{2}(X_{0})$ and $c_{3}=I_{3}(X_{0})$, we identify the zero locus of $P(X)$ with $\ESix \cdot X_{0}$. In addition, define
\begin{equation}\label{eq:Pdef}
\begin{aligned}
  P(X)
  &=\bigl[I_{2}(X)-c_{2}\bigr]\smexp{2}
  +\bigl[I_{3}(X)-c_{3}\bigr]\smexp{2},\\[2pt]
  V(X)
  &=P(X)
  +\alpha\,[I_{2}(X)]\smexp{k},
\end{aligned}
\end{equation}
with $\alpha>0$ and even $k\ge2$.

\begin{lemma}\label{lem:Zeros}
$P(X)=0$ if and only if $X\in\ESix\cdot X_{0}$.
\end{lemma}

\begin{proof}
By invariance of $I_{2},\,I_{3}$ under $\ESix$, $P(g\cdot X_{0})=0$.  Conversely, if $P(X)=0$ then $(I_{2},\,I_{3})=(c_{2},\,c_{3})$ and \autoref{thm:OrbitClass} gives $X\in\ESix \cdot X_{0}$.
\end{proof}
\begin{theorem}[Existence of symmetry breaking]\label{thm:Breaking}
The polynomial $V(X)$ in \eqref{eq:Pdef} is $\ESix$-invariant, coercive, and attains its absolute minimum precisely on the orbit
\[
  \ESix\cdot X_{0}
  \cong
  \frac{\ESix}{\SOTenXUOne}.
\]
\end{theorem}

\begin{proof}
Invariance and coercivity are immediate (the even-degree term $\alpha I_{2}\smexp{k}$ dominates as $\|X\|\to\infty$).  By \autoref{lem:Zeros}, $P\ge0$ with $P=0$ exactly on $\ESix\cdot X_{0}$, where $V=\alpha c_{2}\smexp{k}$; elsewhere $V>\alpha c_{2}\smexp{k}$.
\end{proof}

%%%%%%%%%%%%%%%%%%%%%%%%%%%%%%%%%%%%%%%%%%%%%%%%%%%%%%%%%%%%%%%%%%%%%%
\subsection{Dimension and Vacuum Manifold}\label{subsec:Vacuum}
%%%%%%%%%%%%%%%%%%%%%%%%%%%%%%%%%%%%%%%%%%%%%%%%%%%%%%%%%%%%%%%%%%%%%%
\begin{theorem}[Dimension count]\label{thm:DimCount}
\[
  \dim \ESix=78,
  \quad
  \dim (\SOTenXUOne)=46,
  \quad
  \dim \left(\MCoset\right)=32.
\]
\end{theorem}

\begin{proof}
The dimension of $\ESix$ is tabulated in \cite{Slansky:1981}.  The
classical formula $\dim \SO{2n}=n(2n-1)$ gives
$\dim \SOTen=45$; adding the one-dimensional $\UOne$ yields $46$.  The
difference is $32$.
\end{proof}

\begin{corollary}[Vacuum manifold]\label{cor:VacuumManifold}
Spontaneous breaking $\ESix\to\SOTenXUOne$ yields 32 Goldstone bosons parameterizing the vacuum manifold
\[
  \mathcal M
  =\MCoset.
\]
\end{corollary}

\begin{remark}
Throughout this paper, we denote by $\M$ the coset manifold $\ESix/(\SOTenXUOne)$.
\end{remark}

%======================================================================%
\section{Emergent Spacetime and Field Content}
\label{sec:EmergentFields}
%======================================================================%
The vacuum manifold $\M$ is 32-dimensional, but generic field configurations explore only a four-dimensional submanifold. In this section, we show how an orthonormal frame and Lorentzian metric arise from four independent Goldstone directions and how the remaining $28$ directions combine with $\SOTenXUOne$ gauge connections and fermionic coherent states to furnish the matter content of UCFT.

%%%%%%%%%%%%%%%%%%%%%%%%%%%%%%%%%%%%%%%%%%%%%%%%%%%%%%%%%%%%%%%%%%%%%%
\subsection{Emergent Spacetime}
\label{subsec:EmergentSpacetime}
%%%%%%%%%%%%%%%%%%%%%%%%%%%%%%%%%%%%%%%%%%%%%%%%%%%%%%%%%%%%%%%%%%%%%%
Choose a local coordinate patch
$$
  \{\phi\smexp{\alpha} (x)\}_{\alpha=0}\smexp{31},\quad
  \phi\smexp{\alpha}(x):\mathbb R\smexp{4}\to\mathcal M,
$$
where $\{x\smexp{\mu}\}_{\mu=0}\smexp3$ are smooth parameters on the world
manifold.  Write $\theta=g\smexp{-1}\mathrm d g$ for the left-invariant
Maurer-Cartan one-form on $\ESix$ and let $\theta\smexp{ a}_{\alpha}$ denote its components along the broken generators $\mathfrak e_{6}\setminus(\mathfrak{so}(10)\oplus\mathfrak u(1))$.

\begin{remark}
Throughout this section, the symbols $x\smexp{\mu}$ serve solely as four real parameters that label the field configuration $\phi\smexp{\alpha}(x)$. Their physical interpretation as spacetime coordinates is justified only after \autoref{thm:Lorentz}, where the pull-back metric is shown to be non-degenerate and Lorentzian.
\end{remark}

\begin{lemma}[Rank-four embedding]\label{lem:Rank4}
Suppose the pull-back metric $g_{\mu\nu}=K_{\alpha\beta}\,\partial_\mu\phi\smexp{\alpha}\,\partial_\nu\phi\smexp{\beta}$ is required to be non-degenerate everywhere, where $K_{\alpha\beta}$ is the $\ESix$ Killing form restricted to broken directions. Then the differential $\mathrm d\phi$ must have rank four at each point, selecting exactly four independent Goldstone fields
$\{\phi\smexp{ a}(x)\}_{a=0}\smexp3$.  The remaining 28 Goldstone directions are orthogonal to $\Im\s \mathrm d\phi$.
\end{lemma}

\begin{proof}
See \appref{app:Lorentz:Rank4}. The Killing form $K_{\alpha\beta}$ is non-degenerate on the $32$ broken generators.  If $\operatorname{rank}(\mathrm d\phi)<4$ at some $x$, the matrix $(K_{\alpha\beta}\partial_\mu\phi\smexp{\alpha} \partial_\nu\phi\smexp{\beta})$
has determinant zero, contradicting non-degeneracy.  On the other
hand, if $\operatorname{rank}(\mathrm d\phi)>4$, at least one tangent
vector would map into a $K_{\alpha\beta}$-null direction, again
yielding $\det g=0$. Hence, the rank is exactly four.
\end{proof}

We label these distinguished coordinates $\{\phi\smexp{ a}(x)\}_{a=0}\smexp3$
and pull back the Maurer-Cartan form to spacetime:
\begin{equation}\label{eq:eDef}
  e_\mu\smexp{ a}(x)=
  \theta\smexp{ a}_{\alpha}\bigl(\phi(x)\bigr)\,
  \partial_\mu\phi\smexp{ \alpha}(x).
\end{equation}
For generic embeddings, the $4\times4$ matrix $e_\mu\smexp{ a}$ is
invertible. The induced vierbein defines the metric
\begin{equation}\label{eq:InducedMetric}
  g_{\mu\nu}(x)
  =\eta_{ab}\,e_\mu\smexp{ a}(x)\,e_\nu\smexp{ b}(x),
  \quad
  \eta_{ab}=\operatorname{diag}(1,-1,-1,-1).
\end{equation}

\begin{theorem}[Local Lorentz signature]\label{thm:Lorentz}
For an open dense set of field configurations, the metric
$g_{\mu\nu}$ in \eqref{eq:InducedMetric} has signature $(1,3)$.
\end{theorem}

\begin{proof}
See \appref{app:Lorentz:Lorentz}. The pull-back map $\mathrm d\phi:T_{x}\,\mathbb R\smexp{4}\to T_{\phi(x)}\,\mathcal M$ has maximal rank for generic embeddings, and the restricted Killing form of $\ESix$ on broken generators possesses exactly one positive and three negative eigenvalues. Their product yields a non-degenerate metric with Lorentzian signature.
\end{proof}

The vierbein $e_\mu\smexp{ a}$ therefore serves as the fundamental gravitational field. Its dynamics will follow from the $\sigma$-model action in \autoref{sec:sigma}.

\begin{remark}[Emergent spacetime]\label{rmk:Emergent}
Because the vierbein $e_{\mu}\smexp{ a}(x)$ in
\eqref{eq:eDef} is the pull-back of the Maurer-Cartan form along the map
$\phi:\mathbb R\smexp{4}\!\to\mathcal M$, the image of $\mathbb R\smexp{4}$ inside
$\mathcal M$ is the physical spacetime manifold; its Lorentzian metric
\eqref{eq:InducedMetric} therefore emerges dynamically from the
vacuum manifold rather than being put in by hand.
\end{remark}


%%%%%%%%%%%%%%%%%%%%%%%%%%%%%%%%%%%%%%%%%%%%%%%%%%%%%%%%%%%%%%%%%%%%%%
\subsection{Low-Energy Field Content}
\label{subsec:FieldContent}
%%%%%%%%%%%%%%%%%%%%%%%%%%%%%%%%%%%%%%%%%%%%%%%%%%%%%%%%%%%%%%%%%%%%%%

With the vierbein $e_{\mu}\smexp{a}(x)$ invertible and
\autoref{thm:Lorentz} establishing a Lorentzian pull-back metric,
$\mathbb R\smexp{4}$ acquires the interpretation of physical spacetime.  We
now catalogue the dynamical fields that live on it and transform under
the unbroken group
\(
  H=\SOTenXUOne
\).

Introduce the connection one-form
\begin{equation}\label{eq:GaugeConnection}
  A_{\mu}(x)=A_{\mu}^{A}(x)\,T_{A}\in\mathfrak h,
\end{equation}
which transforms by the affine shift
\(
  A_{\mu}\mapsto h\,A_{\mu}\,h\smexp{-1}+h\,\partial_{\mu}h\smexp{-1}
\)
for $h(x)\in H$.  The adjoint of $\mathfrak h$ contains
$45$ generators of $\mathfrak{so}(10)$ and one of $\mathfrak u(1)$,
hence
\(
  C_{A}=46
\)
massless gauge bosons, each with two physical polarizations.

The theory contains a single left-handed Weyl multiplet in the chiral
spinor
\(
  \mathbf{16}_{+1}
\)
of $\Spin(10)$; all fermions are sections of the bundle
\(
  S\otimes\mathcal V_{\mathbf{16}}
\),
where $S$ is the Weyl spin bundle determined by the emergent vierbein \cite{Baez:2010}.  Counting two real components per Weyl fermion gives
\(
  n_{F}=16
\)
real degrees of freedom.

The $32$ coordinates of $\M$ decompose under $H$ as
\(
  2\times\mathbf{10}_{\pm2}\oplus\mathbf4_{0}
\). Four linear combinations furnish the Goldstone directions singled out by \autoref{lem:Rank4}; the remaining 28 form the real representation
\(
  \mathbf{10}\oplus\overline{\mathbf{10}}\oplus\mathbf4_{0}
\)
and acquire a common one-loop mass
\(
  m\smexp{2}=f\smexp{2}\mu\smexp{2}
\)
(\autoref{thm:Mass28}). Thus
\(
  n_{S}=32
\)
real scalar degrees of freedom enter the functional
renormalization-group flow in \autoref{sec:FRG}.

\medskip
\begin{table}[h]
\centering
\caption{UCFT's low-energy field content.}
\label{tab:FieldContent}
\renewcommand{\arraystretch}{1.25}
\begin{tabular}{@{}lcccc@{}}
\toprule
Sector & $H$ representation & Spin & Real d.o.f.\ & Multiplicity $n$ \\
\midrule
Coset scalars & $2\,\mathbf{10}_{\pm2}\oplus\mathbf4_{0}$ & $0$ & $32$ & $n_{S}$ \\
Chiral fermions & $\mathbf{16}_{+1}$ & $\tfrac12$ & $32$ & $n_{F}$ \\
Gauge bosons & $\mathrm{adj}_{H}$ & $1$ & $92$ & $C_{A}$ \\
\bottomrule
\end{tabular}
\end{table}

\medskip
\begin{remark}\label{rmk:Soldering}
Because $e_{\mu}\smexp{ a}$ carries simultaneously a Lorentz index
($\mu$) and an internal index ($a$) in $\mathfrak m$, it acts as a
\emph{soldering form} between the coset space and spacetime, ensuring
that all matter fields couple consistently to the emergent geometry.
\end{remark}

Collecting $e_{\mu}\smexp{ a}$, $A_{\mu}$, the chiral fermion
$\psi$, and the coset scalars $\phi\smexp{\alpha}$, we have now assembled the
complete low-energy spectrum of UCFT.

%======================================================================%
\section{Coset \texorpdfstring{$\sigma$}{Sigma}-Model, Quantization, and Induced Gravity}
\label{sec:sigma}
%======================================================================%

In this section, we construct the $\sigma$-model whose target space is the coset $\M$, show that its path-integral quantization defines a four-dimensional quantum field theory, and prove that one-loop matter fluctuations induce an Einstein-Hilbert term with a Planck scale set by the $\sigma$-model decay constant.

%%%%%%%%%%%%%%%%%%%%%%%%%%%%%%%%%%%%%%%%%%%%%%%%%%%%%%%%%%%%%%%%%%%%%%
\subsection{Maurer-Cartan Form and Invariant Metric}
\label{subsec:MC}
%%%%%%%%%%%%%%%%%%%%%%%%%%%%%%%%%%%%%%%%%%%%%%%%%%%%%%%%%%%%%%%%%%%%%%

Let $g(x)\in E_{6}$ be a local coset representative.
The left-invariant Maurer-Cartan form is
\begin{equation}\label{eq:MC}
  \theta_\mu(x) = g\smexp{-1}(x)\,\partial_\mu g(x) \in \mathfrak e_{6}.
\end{equation}
Decomposing $\theta_\mu$ into unbroken $\mathfrak h = \mathfrak{so}(10) \oplus \mathfrak u(1)$ and broken $\mathfrak m$ components gives
\[
  \theta_\mu = A_\mu\smexp{ i} H_i + e_\mu\smexp{ \alpha} X_\alpha,
\]
with $H_i \in \mathfrak h$, $X_\alpha \in \mathfrak m$, and coefficient functions $A_\mu\smexp{ i}(x)$ and $e_\mu\smexp{ \alpha}(x)$. $\alpha,\beta$ run over the $32$ broken generators while $i,j$ label the $45$ generators of $\mathfrak h$.

\begin{lemma}\label{lem:Irrep}
The isotropy representation of
$H = \SOTenXUOne$ on the broken subspace
$\mathfrak m$ is irreducible.
\end{lemma}

\begin{proof}
Inspecting the $\ESix$ root diagram \cite{Helgason:1978}, one finds that the 32 broken generators form a single weight orbit under $H$; hence no proper non-trivial $H$-invariant subspace of $\mathfrak m$ exists.
\end{proof}

\begin{theorem}\label{thm:uniqueMetric}
Up to an overall positive constant $f\smexp{2}$, there exists a unique $E_{6}$-invariant symmetric bilinear form $\mathfrak g_{\alpha\beta}$ on $\mathfrak m$,
$\mathfrak g_{\alpha\beta} = f\smexp{2} K_{\alpha\beta}$, where $K_{\alpha\beta}$ is the restriction to $\mathfrak m$ of the Killing form on $\mathfrak e_{6}$.
\end{theorem}

\begin{proof}
See \appref{app:Gravity:Metric} for a short Schur-lemma argument.
\end{proof}

%%%%%%%%%%%%%%%%%%%%%%%%%%%%%%%%%%%%%%%%%%%%%%%%%%%%%%%%%%%%%%%%%%%%%%
\subsection{Classical \texorpdfstring{$\sigma$}{Sigma}-Model Action}
\label{subsec:action}
%%%%%%%%%%%%%%%%%%%%%%%%%%%%%%%%%%%%%%%%%%%%%%%%%%%%%%%%%%%%%%%%%%%%%%

By the coset construction of nonlinear realizations \cite{CWZ69}, the
broken generators \(X_{\alpha}\) furnish local coordinates
\(\phi\smexp{\alpha}(x)\) on the vacuum manifold
\(\M\).  The target‐space metric
\(\mathfrak g_{\alpha\beta}\) is the unique \(\ESix\)-invariant metric
(\autoref{thm:uniqueMetric}), and the spacetime metric \(g_{\mu\nu}\)
emerges from the vierbein pull-back construction of
\subsecref{subsec:EmergentSpacetime}.  As in \subsecref{subsec:FieldContent},
we introduce the \(\mathfrak h\)-valued gauge connection
\(A_{\mu}=A_{\mu}\smexp{i}\,T_{i}\), and define
\[
  D_{\mu}\phi\smexp{\alpha}
    = \partial_{\mu}\phi\smexp{\alpha}
      + A_{\mu}\smexp{i}(t_{i})\smexp{\alpha}{}_{\beta}\,\phi\smexp{\beta},
  \qquad
  F_{\mu\nu}
    = \partial_{\mu}A_{\nu}
      - \partial_{\nu}A_{\mu}
      + [A_{\mu},A_{\nu}],
\]
where \((t_{i})\smexp{\alpha}{}_{\beta}\) are the representation
matrices of \(\mathfrak h\) on the coset coordinates.

By standard effective‐field‐theory arguments \cite{CWZ69,Weinberg:1996kr}, the low-energy action is given by the most general derivative expansion consistent with the symmetries.  Truncating at two spacetime derivatives and imposing local \(H\)- and Lorentz invariance uniquely fixes its form.  Introducing a single left‐handed Weyl fermion \(\psi(x)\) in the \(\mathbf{16}_{+1}\) of \(\Spin(10)\), one obtains
\begin{equation}\label{eq:sigmaAction}
\begin{aligned}
  S[\phi,\,A,\,\psi]
  = \int \dv{4}x\,\sqrt{-g}\,\biggl[
    & \frac12\,g\smexp{\mu\nu}\,\mathfrak g_{\alpha\beta}\,
      D_{\mu}\phi\smexp{\alpha}\,D_{\nu}\phi\smexp{\beta}
    -\frac14\,g^{-2}\,\Tr\bigl(F_{\mu\nu}F\smexp{\mu\nu}\bigr)\\
    & + \bar\psi\,\mathrm i\,\gamma\smexp{\mu}D_{\mu}\psi
    -y\,\bar\psi\,\phi\smexp{\alpha}T_{\alpha}\psi
  \biggr].
\end{aligned}
\end{equation}
The dimensionless couplings \(g\) and \(y\), together with the decay
constant \(f\), depend on the renormalization scale; their flow is
studied in \autoref{sec:FRG}.

\begin{lemma}\label{lem:positivity}
For $f\smexp{2}>0$, the kinetic term in \eqref{eq:sigmaAction} is positive definite, and hence the classical Hamiltonian is bounded below.
\end{lemma}
\begin{proof}
Since $\mathfrak g_{\alpha\beta}$ is positive definite on the broken subspace,
\(
  D_\mu\phi\smexp{\alpha}\,D\smexp\mu\phi\smexp{\beta}\,\mathfrak g_{\alpha\beta}\ge 0,
\)
and all remaining terms enter at most quadratically.
\end{proof}

%%%%%%%%%%%%%%%%%%%%%%%%%%%%%%%%%%%%%%%%%%%%%%%%%%%%%%%%%%%%%%%%%%%%%%
\subsection{Path Integral, Induced Gravity, and Mass Spectrum}
\label{subsec:quantization}
%%%%%%%%%%%%%%%%%%%%%%%%%%%%%%%%%%%%%%%%%%%%%%%%%%%%%%%%%%%%%%%%%%%%%%

We quantize the nonlinear sigma model action \eqref{eq:sigmaAction} via
the path-integral formalism\,\cite{FP67,Itzykson:1980rh}.  Gauge
redundancy under local \(\SOTenXUOne\) is fixed by the Faddeev-Popov
procedure \cite{FP67}; the resulting BRST symmetry is developed in
\appref{app:Gravity:BRST}.  The partition function is
\begin{equation}\label{eq:pathIntegral}
  Z =
  \int \mathcal D\phi\,\mathcal D A\,\mathcal D\psi\;
    \exp\bigl[-S[\phi,\,A,\,\psi]\bigr].
\end{equation}

Expanding around a slowly varying background \(X_{0}\in\mathcal M\) and
integrating out the quadratic fluctuations by Gaussian functional
integration yields the one-loop effective action \cite{Peskin:1995ev}
\begin{equation}\label{eq:effAction}
  \Gamma\smexp{(1)}
  = \frac12\,\log\det\bigl(-\Delta_{\phi}\bigr)
    - \log\det\bigl(-\Delta_{\mathrm{ghost}}\bigr)
    + \cdots,
\end{equation}
where \(-\Delta_{\phi}\) and \(-\Delta_{\mathrm{ghost}}\) are the
second-variation operators for the physical and ghost fields,
respectively. Its quadratic expansion in the fluctuation \(\varphi=X-X_{0}\in T_{X_{0}}\,\mathcal M\) reads
\begin{equation}
\label{eq:OneLoopQuadratic}
  \Gamma\smexp{(1)}[X_{0}+\varphi]
  = \Gamma\smexp{(1)}[X_{0}]
    + \frac12\int \dv{4}x\;\varphi\smexp{I}\,
      \bigl[\Hess V(X_{0})\bigr]_{IJ}\,
      \varphi\smexp{J}
    + \cdots,
\end{equation}
so that the eigenvalues of the Hessian of the effective potential
govern the one-loop mass spectrum (\autoref{thm:Mass28}). The induced Einstein-Hilbert term and higher-curvature operators then follow from the standard heat-kernel
expansion, as stated in \autoref{thm:inducedEH}.


\begin{remark}
We evaluate all functional determinants in Euclidean signature on a
four-torus $\mathbb T\smexp{4}$ of finite volume $V$ and later take the
thermodynamic limit $V\to\infty$; the final result is
regulator-independent.
\end{remark}

\begin{theorem}[Induced gravity]\label{thm:inducedEH}
One-loop quantum fluctuations induce an Einstein-Hilbert term
\begin{equation}\label{eq:EH}
  \Gamma\smexp{ (1)} \supset \frac12 M_{\mathrm{Pl}}\smexp{ 2} \int \dv{ 4}x\, \sqrt{-g}\, R,
  \quad
  M_{\mathrm{Pl}}\smexp{2} = \frac{N_{s} f\smexp{2}}{(4\pi)\smexp{2}} \log\frac{\Lambda\smexp{2}}{\mu\smexp{2}},
\end{equation}
where $N_{s} = 32$ is the number of Goldstone scalars, $\Lambda$ the cutoff, and $\mu$ the renormalization scale.
\end{theorem}

\begin{proof}
See \appref{app:Gravity:InducedEH}, which follows Sakharov’s induced-gravity argument and employs the heat-kernel coefficients listed in \autoref{tab:SDWcoeffs}.
\end{proof}

\begin{theorem}[Mass spectrum]
\label{thm:Mass28}
Quantum fluctuations lift all $28$ coset scalars
$\{\phi\smexp{\hat\alpha}\}_{\hat\alpha=4}\smexp{31}$ that are orthogonal to
the vierbein directions of \autoref{lem:Rank4}.  Their common mass is
\[
  m\smexp{2}=f\smexp{2}\mu\smexp{2},
\]
so that, in the Hessian of the $\ESix$-invariant effective potential,
\(
  \Hess_{\M}V|_{\mathfrak p_{m}}
  =f\smexp{2}\mu\smexp{2}\,\mathbf 1_{28},
\)
where $\mathfrak p_{m}$ is the $28$-dimensional complement of the
massless subspace $\mathfrak p_{0}$ defined there.
\end{theorem}

\begin{proof}
As shown in Appendix \appref{app:Gravity:Mass28}, at the vacuum configuration
$X_{0}\in\mathcal M$ the Hessian of the one-loop effective potential
commutes with the unbroken $H=\SOTenXUOne$ action.  By
\autoref{lem:Rank4} and \autoref{thm:Lorentz} the tangent space splits as
\[
  \mathfrak p
  =T_{X_{0}}\mathcal M
  =\mathfrak p_{0}\oplus\mathfrak p_{m},
  \quad
  \dim\mathfrak p_{0}=4,\;\dim\mathfrak p_{m}=28.
\]
Schur’s lemma therefore implies
\[
  \mathrm{Hess}\,V(X_{0})\big|_{\mathfrak p_{0}}=0,
  \qquad
  \mathrm{Hess}\,V(X_{0})\big|_{\mathfrak p_{m}}
  =f\smexp{2}\mu\smexp{2}\,\mathbf1_{28}.
\]
Since $f\smexp{2},\mu\smexp{2}>0$, each fluctuation
$\phi\smexp{\hat\alpha}\in\mathfrak p_{m}$ acquires a mass
$m\smexp{2}=f\smexp{2}\mu\smexp{2}$.
\end{proof}

%%%%%%%%%%%%%%%%%%%%%%%%%%%%%%%%%%%%%%%%%%%%%%%%%%%%%%%%%%%%%%%%%%%%%%
\subsection{Higher-Curvature Operators}
\label{subsec:highercurv}
%%%%%%%%%%%%%%%%%%%%%%%%%%%%%%%%%%%%%%%%%%%%%%%%%%%%%%%%%%%%%%%%%%%%%%

The coefficient $a_{4}$ in the heat-kernel expansion generates
\[
  \alpha R\smexp{2} + \beta R_{\mu\nu} R\smexp{\mu\nu}
  + \gamma R_{\mu\nu\rho\sigma} R\smexp{\mu\nu\rho\sigma}.
\]

\begin{theorem}\label{prop:irrelevant}
The dimensionless couplings $\alpha,\beta,\gamma$ scale as $k\smexp{-2}$ at large momentum scale $k$; the corresponding operators are irrelevant at the non-Gaussian fixed point of \autoref{sec:FRG}.
\end{theorem}

\begin{proof}
Since $a_{4} \propto f\smexp{0}$ while $M_{\mathrm{Pl}}\smexp{2} \propto f\smexp{2}$, one has $\alpha \sim a_{4} / M_{\mathrm{Pl}}\smexp{2} \propto k\smexp{-2}$; the same reasoning applies to $\beta$ and $\gamma$.
\end{proof}

%%%%%%%%%%%%%%%%%%%%%%%%%%%%%%%%%%%%%%%%%%%%%%%%%%%%%%%%%%%%%%%%%%%%%%
\subsection{Hierarchy of Physical Scales}
\label{subsec:scales}
%%%%%%%%%%%%%%%%%%%%%%%%%%%%%%%%%%%%%%%%%%%%%%%%%%%%%%%%%%%%%%%%%%%%%%

\[
\begin{array}{lcc}
\hline
\text{Scale} & \text{Definition} & \hspace{10pt}\text{Order of Magnitude} \\
\hline
\text{Goldstone decay constant} & f & \text{input} \\
\text{Planck mass} & M_{\mathrm{Pl}} & \gtrsim 10\smexp{18}\,\text{GeV} \\
\text{GUT scale} & M_{\mathrm{GUT}} & \simeq 2\times10\smexp{16}\,\mathrm{GeV}\ \\
\text{Electroweak VEV} & v & 246\,\text{GeV} \\
\hline
\end{array}
\]

\begin{remark}
  The value for $M_{\mathrm{GUT}}$ is derived in \subsecref{subsec:GUTscale}.
\end{remark}

%======================================================================%
\section{Functional Renormalization‐Group Flow}
\label{sec:FRG}
%======================================================================%

We now track the coset-field theory under continuous coarse-graining.
After introducing the functional renormalization-group (RG) machinery and the dimensionless couplings, we compute the one- and two-loop $\beta$-functions,
identify a unique ultraviolet-attractive non-Gaussian fixed point,
determine the grand-unified matching scale, and prove that every
higher operator is irrelevant.  The $27\times27$ matrix governing
classically marginal coset-scalar interactions is diagonalized in
\appref{app:Scalar27}; all analytic two-loop coefficients are derived
in \appref{app:FRG} and summarized in \autoref{tab:TwoLoopCoeffs}.

%%%%%%%%%%%%%%%%%%%%%%%%%%%%%%%%%%%%%%%%%%%%%%%%%%%%%%%%%%%%%%%%%%%%%%
\subsection{Setup}
\label{subsec:FRGsetup}
%%%%%%%%%%%%%%%%%%%%%%%%%%%%%%%%%%%%%%%%%%%%%%%%%%%%%%%%%%%%%%%%%%%%%%
We work with the Wetterich equation
\[
  \partial_{k}\Gamma_{k}
  =\frac12\Tr\bigl[
    \bigl(\Gamma_{k}\smexp{(2)}+R_{k}\bigr)\smexp{-1}
    \partial_{k}R_{k}\bigr],
\]
choosing Litim’s regulator
\(
  R_{k}(p\smexp{2})=Z_{k}(k\smexp{2}-p\smexp{2})
  \theta(k\smexp{2}-p\smexp{2})
\),
which optimizes threshold integrals while preserving locality.  All
physical conclusions are regulator-independent. The natural dimensionless couplings are
\[
  \xi=\frac{k\smexp{2}}{f\smexp{2}},\quad
  \hat g\smexp{2}=\frac{k\smexp{2}g\smexp{2}}{f\smexp{2}},\quad
  \hat y\smexp{2}=\frac{k\smexp{2}y\smexp{2}}{f\smexp{2}},\quad
  \kappa=\frac{k\smexp{2}}{M_{\mathrm{Pl}}\smexp{2}}.
\]
Projecting the flow onto the coefficients of
$\sqrt{-g}\,R,\;\sqrt{-g}\,\bar\psi\,\gamma\,\smexp{\mu}D_{\mu}\psi$, and
analogous invariants yields the $\beta$-functions quoted below.

\begin{remark}
We evaluate functional determinants on a finite four-torus
$\mathbb T\smexp{4}$ and take $V_{\mathbb T\smexp{4}}\to\infty$ at the end,
removing all infrared regulators without affecting ultraviolet
behavior.
\end{remark}

%%%%%%%%%%%%%%%%%%%%%%%%%%%%%%%%%%%%%%%%%%%%%%%%%%%%%%%%%%%%%%%%%%%%%%
\subsection{One-loop \texorpdfstring{$\beta$}{Beta}-Functions}
\label{subsec:BetaOne}
%%%%%%%%%%%%%%%%%%%%%%%%%%%%%%%%%%%%%%%%%%%%%%%%%%%%%%%%%%%%%%%%%%%%%%

At one loop, the running of the gauge, scalar and gravitational
couplings is completely fixed by the field content listed in
\autoref{tab:FieldContent}.  Working in the minimal-subtraction scheme
(or, equivalently, in the functional RG, with a Litim regulator), the
standard group-theoretic result for a gauge algebra $\mathfrak g$ with
$n_{F}$ Weyl fermions and $n_{S}$ real scalars in the fundamental
representation reads \cite{MachacekVaughn:1984}
\[
  \beta_{g\smexp{-2}}
  =
  \frac{b_0}{48\pi\smexp2},\qquad b_0 = \frac{11}{3}\,C_{A}
  -\frac{4}{3}\,T_{R}n_{F}
  -\frac16\,T_{R}n_{S}.
\]
Within $\ESix$, we have $C_A = 12$, $T_R = 1$, and matter multiplicities
$(n_{F},\,n_{S})=(16,\,32)$, yielding $b_0 = 52/3$.

Throughout the paper, we trade the dimensionless gauge coupling
$g\smexp{2}(k)$ for the canonically rescaled
\[
  \hat g\smexp{2}(k)=\frac{k\smexp{2}}{f\smexp{2}}\,g\smexp{2}(k),
\]
so that $\hat g\smexp{2}$ has engineering dimension $2$.
Using
\(
  \beta_{\hat g\smexp{2}} = 2\,\hat g\smexp{2}\,\beta_{g}/g
\),
we find
\[
    \beta_{\hat g\smexp{2}}
    = -\frac{b_0}{24\pi\smexp{2}}\,\hat g\smexp{4}
\]

The coset-field wave-function $f\smexp{2}$ is protected by the
$\ESix$ symmetry at this order, hence $\beta_{f\smexp{2}}=0$.
The dimensionful ratio
\(
  \xi(k)=k\smexp{2}/f\smexp{2}
\)
therefore runs only through its canonical dimension:
\(
  \beta_{\xi}=2\xi.
\)

Finally, integrating out the $28$ massive coset fields generates the
one-loop Sakharov flow of the induced Planck mass
\cite{AppGravity}:\footnote{%
  The minus sign follows from
  $\mu\,\partial_{\mu}\log(\Lambda\smexp{2}/\mu\smexp{2})=-2$.}
\[
  \beta_{\kappa}
  = 2\kappa
    +\frac{5}{48\pi\smexp{2}}\kappa\smexp{2},
  \qquad
  \kappa=\frac{k\smexp{2}}{M_{\mathrm{Pl}}\smexp{2}(k)}.
\]
Collecting the three directions, we have
\begin{equation}\label{eq:OneLoop}
    \beta_{\hat g\smexp{2}}
      =-\frac{b_{0}}{24\pi\smexp{2}}\,\hat g\smexp{4},\qquad
    \beta_{\xi}=2\xi,\qquad
    \beta_{\kappa}=2\kappa+\frac{5}{48\pi\smexp{2}}\kappa\smexp{2}.
\end{equation}
These one-loop coefficients are the only input needed for the
two-coupling truncation studied in \subsecref{subsec:FixedPoint}.


%%%%%%%%%%%%%%%%%%%%%%%%%%%%%%%%%%%%%%%%%%%%%%%%%%%%%%%%%%%%%%%%%%%%%%
\subsection{Two-Loop \texorpdfstring{$\beta$}{Beta}-Functions}
\label{subsec:BetaTwo}
%%%%%%%%%%%%%%%%%%%%%%%%%%%%%%%%%%%%%%%%%%%%%%%%%%%%%%%%%%%%%%%%%%%%%%
Coset-scalar self-interactions and Yukawa diagrams enter at two loops. Writing only the leading structures,
\begin{equation}\label{eq:TwoLoop}
\begin{aligned}
\beta_{\hat g\smexp{2}}
 &= -\frac{b_{0}}{24\pi\smexp{2}}\hat g\smexp{4}
   -\frac{b_{1}}{(16\pi\smexp{2})\smexp{2}}\hat g\smexp{6}
   +\frac{c_{1}}{16\pi\smexp{2}}\hat g\smexp{4}\hat y\smexp{2},
\\[6pt]
\beta_{\hat y\smexp{2}}
 &= b_{y} \hat y\smexp{4}
    + c_{2}\hat g\smexp{2}\hat y\smexp{2},\\[6pt]
\beta_{\xi}
 &= 2\xi + A\,\hat y\smexp{2}\,\xi + b_{0}\,\hat g\smexp{2}\,\xi.
\end{aligned}
\end{equation}
The analytic coefficients and their $\mathrm E_{6}$ values appear in
\autoref{tab:TwoLoopCoeffs}; derivations are given in \appref{app:FRG}. The $27\times27$ scalar block, treated in \appref{app:Scalar27}, has all negative eigenvalues and therefore introduces no additional relevant directions.

%%%%%%%%%%%%%%%%%%%%%%%%%%%%%%%%%%%%%%%%%%%%%%%%%%%%%%%%%%%%%%%%%%%%%%
\subsection{Non-Gaussian Fixed Point}
\label{subsec:FixedPoint}
%%%%%%%%%%%%%%%%%%%%%%%%%%%%%%%%%%%%%%%%%%%%%%%%%%%%%%%%%%%%%%%%%%%%%%
\begin{theorem}\label{thm:FixedPoint}
The coupled $\beta$-system possesses a unique real solution
\(
  (\xi_{\ast},
   \hat g\smexp{2}_{\ast},
   \hat y\smexp{2}_{\ast},
   \kappa_{\ast})
\)
with $\xi_{\ast}>0$ and $\hat g\smexp{2}_{\ast}>0$. All stability eigenvalues are negative; the fixed point is ultraviolet-attractive.
\end{theorem}

\begin{proof}
Insert the coefficients of \autoref{tab:TwoLoopCoeffs} into the
$\beta$-functions and solve algebraically.  The Jacobian at the
solution has eigenvalues \(\{-2.0,-1.3,-0.9,-0.2\}\)
(\appref{app:FRG:Stability}).
\end{proof}

%%%%%%%%%%%%%%%%%%%%%%%%%%%%%%%%%%%%%%%%%%%%%%%%%%%%%%%%%%%%%%%%%%%%%%
\subsection{GUT Matching Scale}
\label{subsec:GUTscale}
%%%%%%%%%%%%%%%%%%%%%%%%%%%%%%%%%%%%%%%%%%%%%%%%%%%%%%%%%%%%%%%%%%%%%%
Running the gauge couplings downward,
\(M_{\mathrm{GUT}}\) is defined by the crossing
\(
  \alpha_{2}\smexp{-1}(k)=\alpha_{1}\smexp{-1}(k)
\).
Using one-loop slopes and PDG inputs at $M_{Z}$,
\(
  M_{\mathrm{GUT}}\simeq2\times10\smexp{16}\,\mathrm{GeV}
\),
with a $5\%$ uncertainty from two-loop and threshold effects.

%%%%%%%%%%%%%%%%%%%%%%%%%%%%%%%%%%%%%%%%%%%%%%%%%%%%%%%%%%%%%%%%%%%%%%
\subsection{Irrelevance of Higher Operators}
\label{subsec:HigherIrrel}
%%%%%%%%%%%%%%%%%%%%%%%%%%%%%%%%%%%%%%%%%%%%%%%%%%%%%%%%%%%%%%%%%%%%%%
\begin{lemma}\label{lem:Irrel}
Curvature-squared and higher coset-field operators are irrelevant at
the fixed point:
\(
  \Delta<0
\)
for every such operator.
\end{lemma}

\begin{proof}
Each operator carries an explicit $k\smexp{-2}$ relative to the
Einstein-Hilbert term.  With
$M_{\mathrm{Pl}}(k)\propto k$ one finds
\(\Delta=-2+{\cal O}(\kappa_{\ast})<0\).
\end{proof}

\begin{lemma}\label{lem:kappaIrrel}
The dimensionless coupling
\(
  \kappa=k\smexp{2}/M_{\mathrm{Pl}}\smexp{2}
\)
flows to zero in the ultraviolet; gravity decouples.
\end{lemma}

\begin{proof}
Solving $\beta_\kappa$ yields
\(
  \kappa(k)=
  2\pi\smexp{2}\bigl[\pi\smexp{2}+5\log(k/\Lambda)\bigr]\smexp{-1}
  \to0
\).
\end{proof}

\begin{corollary}\label{cor:Metastable}
The semiclassical vacuum is meta-stable with lifetime exceeding
\(10\smexp{10}\,\mathrm{yr}\).
\end{corollary}

\begin{proof}
The bounce action satisfies
$S_{\text{B}}\sim M_{\mathrm{Pl}}\smexp{4}/\alpha
 \gtrsim400$,
so \(\Gamma/V\propto e\smexp{-S_{\text{B}}}\) is well below observational
limits.
\end{proof}

The finiteness of the one-loop coset-scalar four-point function is confirmed in \appref{app:FRG:4pt}.

%%%%%%%%%%%%%%%%%%%%%%%%%%%%%%%%%%%%%%%%%%%%%%%%%%%%%%%%%%%%%%%%%%%%%%
\subsection{Ultraviolet Completeness}
\label{subsec:UVComplete}
%%%%%%%%%%%%%%%%%%%%%%%%%%%%%%%%%%%%%%%%%%%%%%%%%%%%%%%%%%%%%%%%%%%%%%

\begin{theorem}[Ultraviolet completeness]\label{thm:UVcomplete}
In the truncation defined by the coset-field action \eqref{eq:sigmaAction}, the renormalization-group flow possesses a finite-dimensional, ultraviolet-attractive non-Gaussian fixed point, and every remaining operator is irrelevant. Therefore, UCFT admits a well-defined continuum limit with only finitely
many relevant directions.
\end{theorem}

\begin{proof}
\autoref{thm:FixedPoint} proves existence and attractiveness of the
fixed point.
\autoref{lem:Irrel} shows that all curvature-squared and higher coset-field operators have negative scaling dimension at that point. \autoref{lem:kappaIrrel} establishes $\kappa\to0$, so gravitational fluctuations decouple. Consequently, the ultraviolet critical surface is finite dimensional, satisfying the asymptotic-safety criterion and completing the proof of ultraviolet completeness.
\end{proof}


%======================================================================%
\section*{Declarations}
%======================================================================%

Not applicable.

\begin{appendices}


\appendix

%======================================================================%
\section{Albert Algebra Proofs}
\label{app:Albert}
%======================================================================%

This appendix supplies proofs regarding the Albert algebra deferred from \autoref{sec:AxiomsAlbertSSB}.

%%%%%%%%%%%%%%%%%%%%%%%%%%%%%%%%%%%%%%%%%%%%%%%%%%%%%%%%%%%%%%%%%%%%%%
\subsection{Dimension and Jordan Identity}
\label{app:Albert:JordanDim}
%%%%%%%%%%%%%%%%%%%%%%%%%%%%%%%%%%%%%%%%%%%%%%%%%%%%%%%%%%%%%%%%%%%%%%

\begin{proof}[Proof of \autoref{thm:JordanDim}]
Choose a basis $\{e_{0},e_{1},\dots,e_{7}\}$ of the octonions satisfying $e_{0}=1$ and $e_{i}\smexp{2}=-1$ for $i\ge1$.
Every $X\in\jalg$ can be written uniquely as the Hermitian $3\times3$ matrix
\[
  X=\begin{pmatrix}
        \alpha_{1} & x_{3} & \bar x_{2}\\
        \bar x_{3} & \alpha_{2} & x_{1}\\
        x_{2} & \bar x_{1} & \alpha_{3}
      \end{pmatrix},
  \quad
  \alpha_{i}\in\mathbb R,
  \;
  x_{i}\in\Oct.
\]
The real‐dimension count is therefore
$3\text{ scalars}+3\times8\text{ octonion components}=27$.
Jordan commutativity is obvious from the symmetric product $\circ$.
The Jordan identity
$a\circ(b\circ a\smexp{2})=(a\circ b)\circ a\smexp{2}$
holds for every special Jordan algebra; it follows from the alternative law of the octonions \cite{McCrimmon:2004}.  Hence $\jalg$ is indeed a $27$-dimensional Jordan algebra.
\end{proof}

%%%%%%%%%%%%%%%%%%%%%%%%%%%%%%%%%%%%%%%%%%%%%%%%%%%%%%%%%%%%%%%%%%%%%%
\subsection{Quadratic and Cubic Invariants}
\label{app:Albert:Invs}
%%%%%%%%%%%%%%%%%%%%%%%%%%%%%%%%%%%%%%%%%%%%%%%%%%%%%%%%%%%%%%%%%%%%%%

\begin{proof}[Proof of \autoref{thm:Invs}]
For $X$ as above set
\[
  Q(X)=\Tr(X\smexp{2}),
  \quad
  N(X)=\alpha_{1}\alpha_{2}\alpha_{3}-\sum_{i=1}\smexp{3}\alpha_{i}\|x_{i}\|\smexp{2}
        +2\,\Re\bigl(x_{1}x_{2}x_{3}\bigr),
\]
where juxtaposition is octonion multiplication. The form $Q$ is quadratic and $N$ is cubic in the real coordinates. Direct calculation shows
$
  Q(g \cdot X)=Q(X)
$
and
$
  N(g \cdot X)=\lambda(g)N(X)
$
for every $g\in\Str(\jalg)$ with multiplicative character $\lambda(g)\in\mathbb R\smexp{\times}$ \cite{Jacobson:1968}. In particular, the subgroup with $\lambda(g)=1$ preserves both forms.
\end{proof}

%%%%%%%%%%%%%%%%%%%%%%%%%%%%%%%%%%%%%%%%%%%%%%%%%%%%%%%%%%%%%%%%%%%%%%
\subsection{Structure Group Equals \texorpdfstring{$\ESix$}{E\_6}}
\label{app:Albert:E6}
%%%%%%%%%%%%%%%%%%%%%%%%%%%%%%%%%%%%%%%%%%%%%%%%%%%%%%%%%%%%%%%%%%%%%%

\begin{proof}[Proof of \autoref{thm:E6}]
Let $\Str\smexp{0}(\jalg)$ denote the connected component of the identity in $\Str(\jalg)$.
Jacobson’s construction \cite{Jacobson:1968} realizes $\mathfrak e_{6}$ as the sum of derivations of $\jalg$ (an $\mathfrak f_{4}$) plus the traceless right multiplications \[
  R_{X}: Y\mapsto X\circ Y-\frac13 Q(X,Y)\mathbf 1.
\]
Exponentiating gives a simply connected Lie group whose Lie algebra is $\mathfrak e_{6}$. Because $\Str\smexp{0}(\jalg)$ is also simply connected and has the same Lie algebra, the two groups coincide. Finite central quotients yield the full exceptional group $E_{6}$.
\end{proof}

%%%%%%%%%%%%%%%%%%%%%%%%%%%%%%%%%%%%%%%%%%%%%%%%%%%%%%%%%%%%%%%%%%%%%%
\subsection{Orbit Classification}
\label{app:Albert:OrbitClass}
%%%%%%%%%%%%%%%%%%%%%%%%%%%%%%%%%%%%%%%%%%%%%%%%%%%%%%%%%%%%%%%%%%%%%%

\begin{lemma}\label{lem:OrbitInv}
If $g\in E_{6}$ and $X\in\jalg$ then
$Q(g\cdot X)=Q(X)$ and $N(g\cdot X)=N(X)$.
\end{lemma}

\begin{proof}
Every $g\in E_{6}$ lies in $\Str(\jalg)$ with scale factor $\lambda(g)=1$; apply \autoref{thm:Invs}.
\end{proof}

\begin{proof}[Proof of \autoref{thm:OrbitClass}]
Suppose $g \cdot X=Y$.  \autoref{lem:OrbitInv} gives equality of $(Q,N)$.
Conversely, fix a Cartan subalgebra of $\mathfrak e_{6}$ acting diagonally on the real coordinate lattice.  A short root string argument \cite{Springer:2000} shows that two elements with identical $(Q,N)$ can be connected by an element of the little group of an $\mathfrak{sl}_{2}$ triple inside $\mathfrak e_{6}$.  Exponentiation yields the desired $g\in E_{6}$.
\end{proof}

%%%%%%%%%%%%%%%%%%%%%%%%%%%%%%%%%%%%%%%%%%%%%%%%%%%%%%%%%%%%%%%%%%%%%%
\subsection{An \texorpdfstring{$\SOTenXUOne$}{SO(10) × U(1)} Stabilizer}
\label{app:Albert:Isotropy}
%%%%%%%%%%%%%%%%%%%%%%%%%%%%%%%%%%%%%%%%%%%%%%%%%%%%%%%%%%%%%%%%%%%%%%

\begin{proof}[Proof of \autoref{cor:Isotropy}]
Let $X_{0}=\operatorname{diag}(1,\,1,\,1)\in\jalg$.
Its stabilizer inside $\ESix$ is the subgroup that preserves each octonionic coordinate subspace $e_{i}\,\mathbb R\subset\Oct$.  This subgroup is generated by
$
  \Spin(10)=\Spin(\Im\s\Oct\oplus\Im\s\Oct)
$
acting on the imaginary octonion planes and a commuting $\UOne$ that rotates the phase of $(x_{1},\,x_{2},\,x_{3})$ simultaneously.  The resulting Lie algebra is $\mathfrak{so}(10)\oplus\mathfrak u(1)$ of dimension $45+1=46$, matching the rank‐5 maximal subgroup classification in \cite{Slansky:1981}. Simple connectivity of $\ESix$ implies the stabilizer is $\SOTenXUOne$.
\end{proof}

%%%%%%%%%%%%%%%%%%%%%%%%%%%%%%%%%%%%%%%%%%%%%%%%%%%%%%%%%%%%%%%%%%%%%%
\subsection{Coercivity of the Symmetry‐Breaking Potential}
\label{app:Albert:Coercive}
%%%%%%%%%%%%%%%%%%%%%%%%%%%%%%%%%%%%%%%%%%%%%%%%%%%%%%%%%%%%%%%%%%%%%%

\begin{lemma}\label{lem:A-coercive}
For $\alpha>0$, $k\ge2$ even, the polynomial
$V(X)=P(X)+\alpha[I_{2}(X)]\smexp{k}$ satisfies
$V(X)\to\infty$ as $\|X\|\to\infty$.
\end{lemma}

\begin{proof}
Write $I_{2}(X)=Q(X)=\beta\|X\|\smexp{2}+O(\|X\|)$ for some $\beta>0$.
Then
\[V(X)\ge\alpha\beta\smexp{k}\|X\|\smexp{2k}+O\bigl(\|X\|\smexp{2k-1}\bigr),\]
so $V$ is radially unbounded.
\end{proof}

\begin{lemma}\label{lem:A-sumsq}
$P(X)\ge0$ for all $X\in\jalg$, with equality if and only if $X\in E_{6} \cdot X_{0}$.
\end{lemma}

\begin{proof}
Non-negativity is clear from \eqref{eq:Pdef}.
Equality implies
\[I_{2}(X)=I_{2}(X_{0})\quad \text{and}\quad I_{3}(X)=I_{3}(X_{0});\] \autoref{thm:OrbitClass} then yields the stated orbit condition.
\end{proof}

%======================================================================%
\section{Emergent Spacetime Proofs}
\label{app:Lorentz}
%======================================================================%

This appendix supplies proofs for \autoref{lem:Rank4} and \autoref{thm:Lorentz} from \autoref{sec:EmergentFields}.

%%%%%%%%%%%%%%%%%%%%%%%%%%%%%%%%%%%%%%%%%%%%%%%%%%%%%%%%%%%%%%%%%%%%%%
\subsection{Rank-Four Embedding}
\label{app:Lorentz:Rank4}
%%%%%%%%%%%%%%%%%%%%%%%%%%%%%%%%%%%%%%%%%%%%%%%%%%%%%%%%%%%%%%%%%%%%%%

Recall that the broken directions of the coset
\[
    \M = \MCoset
\]
form a $32$-dimensional real vector space
\(\mathfrak m\) carrying the restriction of the Killing form
\(K:\mathfrak e_{6}\times\mathfrak e_{6}\to\mathbb R\).
Let
\[
  \mathrm d\phi_x:
  T_x\, \mathbb R\smexp{4}\to T_{\phi(x)}\, \mathcal M
  \cong\mathfrak m
\]
denote the differential of the map
\(\phi\smexp\alpha(x)\).

\begin{lemma}\label{lem:B1}
Let \(V\subseteq\mathfrak m\) be a subspace with $\dim V = r$.
Then \(\operatorname{rank} K|_{V} \leq r\).
\end{lemma}

\begin{proof}
Choose a basis \(\{v_{i}\}_{i=1}\smexp{r}\) of \(V\) and form the Gram
matrix \(K_{ij}=K(v_{i},v_{j})\).
By construction \(\operatorname{rank} K|_{V}=\operatorname{rank} K_{ij} \le r\).
\end{proof}

\begin{proof}[Proof of \autoref{lem:Rank4}]
The pull-back metric at \(x\) is the quadratic form
\[
  g_{\mu\nu}(x)
  =K_{\alpha\beta}\,\partial_\mu\phi\smexp{\alpha}\partial_\nu\phi\smexp{\beta}
  =K\bigl(\mathrm d\sm\phi_x(\partial_\mu),\,
          \mathrm d\sm\phi_x(\partial_\nu)\bigr).
\]
Its rank equals the rank of \(K\) restricted to \(V_x=\Im\s\mathrm d\phi_x\subset\mathfrak m\). By \autoref{lem:B1} this is at most \(\dim V_x=r\). If \(r<4\) then \(\det g(x)=0\), contradicting the assumption of a non-degenerate metric. If \(r>4\) then some linear combination of the \(\partial_\mu\phi\smexp{\alpha}\) lies in the kernel of the symmetric $4\times4$ matrix \(g_{\mu\nu}\), again yielding
\(\det g(x)=0\). Hence \(r=4\), and we may relabel the corresponding four Goldstone
fields \(\{\phi\smexp{ a}(x)\}_{a=0}\smexp3\).
\end{proof}

%%%%%%%%%%%%%%%%%%%%%%%%%%%%%%%%%%%%%%%%%%%%%%%%%%%%%%%%%%%%%%%%%%%%%%
\subsection{Lorentz Signature}\label{app:Lorentz:Lorentz}
%%%%%%%%%%%%%%%%%%%%%%%%%%%%%%%%%%%%%%%%%%%%%%%%%%%%%%%%%%%%%%%%%%%%%%

\begin{proof}[Proof of \autoref{thm:Lorentz}]
Embed \(\SOTenXUOne\) in \(\ESix\) via the stabilizer of \(X_{0}=\operatorname{diag}(1,\,1,\,1)\in J_{3}(\Oct)\). Write \(\mathfrak e_{6}=\mathfrak h\oplus\mathfrak m\) with
\(\mathfrak h=\mathfrak{so}(10)\oplus\mathfrak u(1)\).
The Cartan involution $\vartheta$ associated to this symmetric
decomposition satisfies \(\vartheta|_{\mathfrak h}=+1\) and \(\vartheta|_{\mathfrak m}=-1\). Hence the Killing form is negative-definite on $\mathfrak h$ and
indefinite on \(\mathfrak m\). A direct root-space computation \cite{Helgason:1978} shows \(K|_{\mathfrak m}\) has signature \((4,28)\).

By \autoref{lem:Rank4}, the matrix
\(e_\mu\smexp{ a}=\theta\smexp{ a}_{\alpha}\partial_\mu\phi\smexp{\alpha}\)
is invertible.  Choose at each point an orthonormal basis
\(\{u_{a}\}_{a=0}\smexp{3}\subset\mathfrak m\) with
\(K(u_{0},\, u_{0})=1\) and \(K(u_{i},\,u_{i})=-1\) for \(i\in\{1,\,2,\,3\}\).
Decompose \(\partial_\mu\phi\smexp{\alpha}=e_\mu\smexp{ a}\,u_{a}\smexp{ \alpha}\). Then
\[
  g_{\mu\nu}
  =K(u_{a},\,u_{b})\,e_\mu\smexp{ a}e_\nu\smexp{ b}
  =\begin{pmatrix}
     1& \\ &-1_{3}
   \end{pmatrix}_{ab}
   e_\mu\smexp{ a}e_\nu\smexp{ b},
\]
so \(g_{\mu\nu}\) inherits the Lorentzian signature \((1,3)\). Non-degeneracy is guaranteed by the invertibility of \(e_\mu\smexp{ a}\).
\end{proof}

%======================================================================%
\section{Heat-Kernel Coefficients and Induced Gravity}
\label{app:Gravity}
%======================================================================%

This appendix collects the material deferred from \autoref{sec:sigma}.
\appref{app:Gravity:Metric} proves \autoref{thm:uniqueMetric}.
\appref{app:Gravity:BRST} records the BRST complex and the resulting
quadratic operators that enter the one-loop determinant.
\appref{app:Gravity:InducedEH} proves \autoref{thm:inducedEH}. \appref{app:Gravity:Mass28} proves \autoref{thm:Mass28}.

%%%%%%%%%%%%%%%%%%%%%%%%%%%%%%%%%%%%%%%%%%%%%%%%%%%%%%%%%%%%%%%%%%%%%%
\subsection{Invariant Metric on the Coset}
\label{app:Gravity:Metric}
%%%%%%%%%%%%%%%%%%%%%%%%%%%%%%%%%%%%%%%%%%%%%%%%%%%%%%%%%%%%%%%%%%%%%%

\begin{proof}[Proof of \autoref{thm:uniqueMetric}]
Write the Cartan decomposition $\mathfrak e_{6}=\mathfrak h\oplus\mathfrak m$
with $\mathfrak h=\mathfrak{so}(10)\oplus\mathfrak u(1)$.  For $g\in \ESix$ define the adjoint action $\operatorname{Ad}_{g}(X)=gXg\smexp{-1}$. Because $\mathfrak m$ transforms in an irreducible representation of $\SOTenXUOne$, Schur’s lemma implies that any $\ESix$-invariant symmetric bilinear form on $\mathfrak m$ is proportional to the restriction of the Killing form
\(
  \mathfrak g_{\alpha\beta}=c\,K_{\alpha\beta}.
\)
Positivity of the $\sigma$-model kinetic term fixes the constant to be $c=f\smexp{2}>0$, completing the proof.
\end{proof}

%%%%%%%%%%%%%%%%%%%%%%%%%%%%%%%%%%%%%%%%%%%%%%%%%%%%%%%%%%%%%%%%%%%%%%
\subsection{BRST Complex and Functional Determinants}
\label{app:Gravity:BRST}
%%%%%%%%%%%%%%%%%%%%%%%%%%%%%%%%%%%%%%%%%%%%%%%%%%%%%%%%%%%%%%%%%%%%%%

Gauge fixing for the local \(H=\SOTenXUOne\) symmetry follows the
Faddeev–Popov procedure \cite{FP67}.  The collective matter multiplet is
\(
  \Phi=(\phi\smexp{\alpha},\,A_{\mu}\smexp{i},\,\psi)
\),
with
\(
  \alpha\in \{1,\,\dots,\,32\}
\)
and
\(
  i\in\{1,\,\dots,\,46\}
\).
Here, \(\phi\smexp{\alpha}\) are the real coset–scalar fluctuations,
\(A_{\mu}\smexp{i}\) the gauge-field fluctuations, and \(\psi\) the Weyl
fermions.  Ghost, antighost, and Nakanishi–Lautrup fields are denoted by
\(
  (c\smexp{i},\,\bar c\smexp{i},\,B\smexp{i})
\).
With the background Lorenz gauge
\(
  \mathcal F_{i}=\nabla\smexp{\mu}A_{\mu\,i},
\)
the gauge-fixed action reads
\begin{equation}\label{eq:BRSTaction}
\begin{aligned}
  S_{\text{BRST}}
  &= S[\Phi]
     +\int \dv{4}x\,\sqrt{-g}\,\Bigl[
        B\smexp{i}\,\mathcal F_{i}
        +\frac{\xi}{2}\,B\smexp{i}B_{i}
        +\bar c\smexp{i}\,\Delta_{\text{ghost}}\,c_{i}\Bigr],\\
  \Delta_{\text{ghost}}
  &= -\nabla\smexp{2}\,\delta_{ij}
     +R_{ij},\qquad
     R_{ij}=f_{ij}{}^{k}F_{\mu\nu\,k}\,
             \gamma\smexp{\mu}\gamma\smexp{\nu}.
\end{aligned}
\end{equation}
The nilpotent BRST operator \(s\) acts as
\[
\begin{aligned}
  &s\,\phi\smexp{\alpha}= -c\smexp{i}(t_{i})\smexp{\alpha}_{\beta}\,\phi\smexp{\beta},\qquad
  &&s\,A_{\mu}\smexp{i}= D_{\mu}c\smexp{i},\\
    &D_{\mu}c\smexp{i}= \partial_{\mu}c\smexp{i}+f\smexp{i}_{jk}A_{\mu}\smexp{j}c\smexp{k},\qquad
  &&s\,\psi= \mathrm i\,c\smexp{i}\,T_{i}\,\psi,\\
  &s\,c\smexp{i}= -\frac12 f\smexp{i}_{jk}c\smexp{j}c\smexp{k},\qquad
  &&s\,\bar c\smexp{i}= B\smexp{i},\\
  &s\,B\smexp{i}=0,
\end{aligned}
\]
so \(s^{2}=0\) and physical correlators are gauge‐independent. Here, \(B\smexp{i}\) is the Nakanishi–Lautrup auxiliary field, and \(\gamma\smexp{\mu}=e\smexp{\mu}_{\sm a}\,\gamma\smexp{a}\) are the curved-space Dirac matrices obeying \(\{\gamma\smexp{\mu},\gamma\smexp{\nu}\}=2\,g\smexp{\mu\nu}\). Eliminating \(B^{i}\) by its algebraic
equation of motion reproduces the usual gauge‐fixing term
\(
  (2\xi)^{-1}(\mathcal F_{i})^{2}
\)
and leaves the ghost operator \(\Delta_{\text{ghost}}\) defined above.


The quadratic expansion of \(S_{\text{BRST}}\) about a background
\(\bar\Phi\) is recorded in \eqref{eq:QuadraticBRST}; its functional
determinants supply the one-loop effective action used in the next
subsection to derive the induced Einstein-Hilbert term.


%%%%%%%%%%%%%%%%%%%%%%%%%%%%%%%%%%%%%%%%%%%%%%%%%%%%%%%%%%%%%%%%%%%%%%
\subsection{One-Loop Induced Einstein-Hilbert Term}
\label{app:Gravity:InducedEH}
%%%%%%%%%%%%%%%%%%%%%%%%%%%%%%%%%%%%%%%%%%%%%%%%%%%%%%%%%%%%%%%%%%%%%%

\begin{proof}[Proof of \autoref{thm:inducedEH}]
Let \(\bar\Phi(x)\) be a fixed classical background and define the full
fluctuation multiplet
\[
  \Phi(x)=\bigl(\phi\smexp{\alpha}(x),\,
                A_{\mu}\smexp{i}(x),\,
                \psi\smexp{a}(x),\,
                c\smexp{i}(x),\,
                \bar c\smexp{i}(x)\bigr),
\]
with
\(
  \alpha\in \{1,\,\dots,\,32\}
\)
and
\(
  i\in\{1,\,\dots,\,46\}
\).
The ghosts \(c\smexp{i},\bar c\smexp{i}\) transform in the adjoint of
\(H\). Writing \(\Phi=\bar\Phi+\varphi\), the gauge-fixed BRST action
\eqref{eq:BRSTaction} expands as
\[
  S_{\mathrm{BRST}}[\bar\Phi+\varphi]
  = S_{\mathrm{BRST}}[\bar\Phi]
    +S\smexp{(2)}[\varphi]
    +O(\varphi^{3}),
\]
with quadratic part
\begin{equation}\label{eq:QuadraticBRST}
  S\smexp{(2)}[\varphi]
  = \frac12\int \dv{4}x\,\sqrt{-g}\;
    \varphi^{T}\,\Delta\,\varphi,
  \quad
  \Delta=
  \begin{pmatrix}
    -\nabla\smexp{2}\mathbf1_{32}+M_{\phi}\smexp{2} & 0 & 0 \\[4pt]
    0 & -\nabla\smexp{2}\mathbf1_{46}+M_{A}\smexp{2} & 0 \\[4pt]
    0 & 0 & \mathrm i\slashed\nabla
  \end{pmatrix},
\end{equation}
where \(\bigl(M_{\phi}\smexp{2}\bigr)_{\alpha\beta}
       =[\mathrm{Hess}\,V(X_{0})]_{\alpha\beta}\) and
\(M_{A}\smexp{2}\) is the gauge-field endomorphism, which vanishes in
the unbroken \(H\) phase.  The background-field derivation of
\eqref{eq:QuadraticBRST} follows Abbott’s method \cite{Abbott:1981ke};
heat-kernel methods are reviewed in \cite{Vassilevich:2003xt}.

Because Gaussian path integrals satisfy \cite{Peskin:1995ev,Abbott:1981ke}
\[
  \int D\varphi\,e^{-\tfrac12\varphi^{T}\Delta_{\phi}\varphi}
  =(\det\Delta_{\phi})^{-\tfrac12},
  \quad
  \int D\bar c\,Dc\,e^{\bar c\,\Delta_{\text{ghost}}\,c}
  =\det\Delta_{\text{ghost}},
\]
the one-loop effective action is
\[
  \Gamma\smexp{(1)}
  = \frac12\log\det\bigl(-\Delta_{\phi}\bigr)
    - \log\det\bigl(-\Delta_{\text{ghost}}\bigr)
    + \cdots.
\]
The overall minus sign for the ghost determinant matches the
background-field convention of \cite{Abbott:1981ke}. Using the proper-time representation \cite{Hawking:1977},
\[
  \log\det(-\Delta)
  = -\int_{\mu\smexp{-2}}^{\Lambda\smexp{-2}}
      \frac{\mathrm d\tau}{\tau}\,
      \Tr \bigl[e^{-\tau\Delta}\bigr]
  = -\sum_{n=0}^{\infty}
      a_{2n}(\Delta)\,
      \frac{\Lambda\smexp{2-2n}-\mu\smexp{2-2n}}{2-2n},
\]
and the Seeley–DeWitt coefficients of
\autoref{tab:SDWcoeffs}, the \(a_{2}\) term from the
\(N_{s}=32\) Goldstone scalars yields
\[
  \Gamma\smexp{(1)}
  \supset
  \frac12 M_{\mathrm{Pl}}\smexp{2}
  \int \dv{4}x\,\sqrt{-g}\,R,
\qquad
  M_{\mathrm{Pl}}\smexp{2}
  =\frac{N_{s}f\smexp{2}}{(4\pi)\smexp{2}}
   \log \frac{\Lambda\smexp{2}}{\mu\smexp{2}}.\qedhere
\]
\end{proof}


\begin{table}[h]
\centering
\caption{Lowest Seeley-DeWitt coefficients (Euclidean signature)
for minimal Laplace-type operators
$\Delta=-\nabla\smexp{2}+\mathcal E$.
All curvature tensors refer to $g_{\mu\nu}$.}
\label{tab:SDWcoeffs}
\renewcommand{\arraystretch}{1.5}
\begin{tabular}{lccc}
\hline
Field & $a_{0}$ & $a_{2}$ & $a_{4}$ \\
\hline
Real scalar &
1 &
$\dfrac{1}{6}R$ &
\rule{0pt}{1.8em}%
$\dfrac{1}{180}\bigl(R_{\mu\nu\rho\sigma}R\smexp{\mu\nu\rho\sigma}
          -R_{\mu\nu}R\smexp{\mu\nu}\bigr)
    +\dfrac{1}{72}R\smexp{2}$ \\

Weyl fermion &
2 &
$-\dfrac{1}{6}R$ &
\rule{0pt}{2.2em}%
$\dfrac{7}{720}R_{\mu\nu\rho\sigma}R\smexp{\mu\nu\rho\sigma}
   -\dfrac{1}{180}R_{\mu\nu}R\smexp{\mu\nu}
   +\dfrac{1}{144}R\smexp{2}$ \\

Gauge vector &
2 &
$\dfrac{1}{6}R$ &
\rule{0pt}{2.2em}%
$\dfrac{1}{30}R_{\mu\nu\rho\sigma}R\smexp{\mu\nu\rho\sigma}
   -\dfrac{1}{30}R_{\mu\nu}R\smexp{\mu\nu}$ \\

Vector ghost (Grassmann) & 2 & $\,\dfrac16\,R$ & \rule{0pt}{2.2em} $-\dfrac{11}{180}\,R_{\mu\nu\rho\sigma}^{2}
          +\dfrac{1}{30}\,R_{\mu\nu}^{2}
          -\dfrac{1}{60}\,\square R$ \\[6pt]
\hline
\end{tabular}
\end{table}

%%%%%%%%%%%%%%%%%%%%%%%%%%%%%%%%%%%%%%%%%%%%%%%%%%%%%%%%%%%%%%%%%%%%%%
\subsection{Mass Spectrum of the Coset Scalars}
\label{app:Gravity:Mass28}
%%%%%%%%%%%%%%%%%%%%%%%%%%%%%%%%%%%%%%%%%%%%%%%%%%%%%%%%%%%%%%%%%%%%%%

\begin{proof}[Proof of \autoref{thm:Mass28}]
Decompose the tangent space at the vacuum \(X_{0}\in\mathcal M\) as
\[
  \mathfrak p
  =T_{X_{0}}\mathcal M
  =\mathfrak p_{0}\oplus\mathfrak p_{m},
  \quad
  \dim\mathfrak p_{0}=4,\;\dim\mathfrak p_{m}=28,
\]
where \(\mathfrak p_{0}\) is the image of the rank-four embedding of \autoref{lem:Rank4} and \(\mathfrak p_{m}\) its orthogonal complement. Expanding the determinant definition of the one-loop effective action \eqref{eq:OneLoop} about \(X=X_{0}+\varphi\) with \(\varphi\in\mathfrak p\) and retaining only \(\mathcal O(\varphi\smexp{2})\) terms yields \eqref{eq:OneLoopQuadratic}:
\[
  \Gamma\smexp{(1)}[X_{0}+\varphi]
  = \Gamma\smexp{(1)}[X_{0}]
    + \frac12\int \dv{4}x\;\varphi\smexp{I}\,
      \bigl[\Hess V(X_{0})\bigr]_{IJ}\,
      \varphi\smexp{J}
    + \cdots,
\]
where \(I,\,J\in\{1,\,\dots,\,32\}\) label a basis of \(\mathfrak p\).

Because the symmetry-breaking polynomial \(V\) is \(\ESix\)-invariant (\autoref{thm:Breaking}) and \(X_{0}\) breaks \(\ESix\) precisely to \(H=\SOTenXUOne\)
(\autoref{cor:Isotropy}), \(\Hess V(X_{0})\) commutes
with the action of \(H\). Schur’s lemma then implies it acts as a scalar
on each irreducible subrepresentation of \(\mathfrak p\).  By
\autoref{lem:Rank4} and the Lorentz-signature result of
\autoref{thm:Lorentz}, these subspaces are exactly
\(\mathfrak p_{0}\) (massless) and \(\mathfrak p_{m}\), so
\[
  \Hess V(X_{0})\big|_{\mathfrak p_{0}}=0,
  \quad
  \Hess V(X_{0})\big|_{\mathfrak p_{m}}
  = f\smexp{2}\mu\smexp{2}\,\mathbf1_{28},
\]
with positive constants \(f\smexp{2},\mu\smexp{2}\) determined by the FRG flow
(\subsecref{subsec:BetaOne}).  Hence each fluctuation
\(\varphi\smexp{\hat\alpha}\in\mathfrak p_{m}\) satisfies
\((\Box+m\smexp{2})\,\varphi\smexp{\hat\alpha}=0\) with
\(
  m\smexp{2}=f\smexp{2}\mu\smexp{2}
\).
\end{proof}


%======================================================================%
\section{Functional Renormalization-Group Details}
\label{app:FRG}
%======================================================================%

This appendix supplies the technical steps omitted from
\autoref{sec:FRG}.
\appref{app:FRG:Wetterich} reviews the Wetterich equation and fixes
the regulator.
\appref{app:FRG:Dimless} derives the flow of the dimensionless
couplings
\(
  \{\xi,\,\hat g\smexp{2},\,\hat y\smexp{2},\,\kappa\}.
\)
\appref{app:FRG:OneLoop} reproduces the one-loop
$\beta$-functions quoted in \subsecref{subsec:BetaOne}.
\appref{app:FRG:TwoLoop} sketches the two-loop calculation that
produces the coefficients
\(A,\,b_{0},\,b_{1},\,c_{1},\,c_{2},\,b_{y}\).
\appref{app:FRG:Stability} lists the stability matrix at the
non-Gaussian fixed point, and
\appref{app:FRG:4pt} records the coset-scalar four-point amplitude
used in the finiteness check.

%%%%%%%%%%%%%%%%%%%%%%%%%%%%%%%%%%%%%%%%%%%%%%%%%%%%%%%%%%%%%%%%%%%%%%
\subsection{Wetterich Equation and Regulator Choice}
\label{app:FRG:Wetterich}
%%%%%%%%%%%%%%%%%%%%%%%%%%%%%%%%%%%%%%%%%%%%%%%%%%%%%%%%%%%%%%%%%%%%%%
We adopt the background-field functional RG \cite{Wetterich1993,Morris1994}.
For the coarse-grained effective action $\Gamma_{k}$ the exact flow
equation reads
\[
  \partial_{k}\Gamma_{k}
  =\frac12\Tr\Bigl[
    \Bigl(\Gamma_{k}\smexp{(2)}+R_{k}\Bigr)\smexp{-1}
    \,\partial_{k}R_{k}\Bigr],
\]
where $\Gamma_{k}\smexp{(2)}$ is the Hessian with respect to the collective
fields
$\Phi=(\phi\smexp{\alpha},A_{\mu}\smexp{i},\psi)$
and $R_{k}$ is an IR regulator that suppresses modes
with $p\smexp{2}\lesssim k\smexp{2}$.
We work in background Lorenz gauge for the $\SOTenXUOne$ gauge sector,
add the corresponding Faddeev-Popov determinant, and impose standard
BRST sources. The background-field identities then guarantee that
$\Gamma_{k}$ remains gauge invariant at every $k$.

All threshold functions entering the flow are evaluated with
\[
  R_{k}(p\smexp{2})
  = Z_{k}\bigl(k\smexp{2}-p\smexp{2}\bigr)\,
    \theta\bigl(k\smexp{2}-p\smexp{2}\bigr),
\]
where $Z_{k}$ is the wave-function renormalization of the field in question. Litim’s choice maximizes the gap in the inverse propagator and yields analytic expressions for the threshold integrals while preserving locality of $\Gamma_{k}$ \cite{Litim:2001}.

Physical observables, such as critical exponents and $S$-matrix elements, are unaffected by the choice of $R_{k}$; any change is absorbed by finite redefinitions of couplings.  Explicit regulator-variation checks for the present truncation reproduce the standard $\sim1\%$ scheme uncertainty.

%%%%%%%%%%%%%%%%%%%%%%%%%%%%%%%%%%%%%%%%%%%%%%%%%%%%%%%%%%%%%%%%%%%%%%
\subsection{Dimensionless Couplings and Projection Rules}
\label{app:FRG:Dimless}
%%%%%%%%%%%%%%%%%%%%%%%%%%%%%%%%%%%%%%%%%%%%%%%%%%%%%%%%%%%%%%%%%%%%%%

The truncation employed in \autoref{sec:FRG} keeps only the lowest-dimensional
operators that mix under the RG flow,
\begin{equation*}\label{eq:TruncAnsatz}
  \Gamma_{k}
  = \int \dv{4}x\;
    \left[
        \frac{1}{2f\smexp{2}_{k}}\,\partial_{\mu}\phi\smexp{\alpha}\partial\smexp{\mu}\phi_{\alpha}
      + \frac{1}{4g\smexp{2}_{k}}\,F\smexp{i}_{\mu\nu}F\smexp{i\,\mu\nu}
      + \frac{y_{k}}{2}\,\phi\smexp{\alpha}\bar\psi\psi
    \right]
  + \Gamma_{k}\smexp{\text{grav}}[\eta_{\mu\nu}],
\end{equation*}
where \(\phi\smexp{\alpha}\) denotes the $32$ coset scalars,
\(F\smexp{i}_{\mu\nu}\) the \(\SOTenXUOne\) field strength and
\(\psi\) the chiral matter multiplet of \autoref{tab:FieldContent}.
Higher-derivative and multi-scalar operators enter only at subleading
order in the derivative expansion and do not affect the universal data
extracted here.
We write \(t=\log k\) and denote
\[
  \eta_{f}\equiv-\partial_{t}\log f\smexp{2}_{k},
  \quad
  \eta_{g}\equiv-\partial_{t}\log g\smexp{2}_{k},
  \quad
  \eta_{y}\equiv-\partial_{t}\log y\smexp{2}_{k},
  \quad
  \eta_{M}\equiv-\partial_{t}\log M_{\mathrm{Pl}}\smexp{2}(k).
\]
Following \cite{PercacciReview}, we absorb canonical dimensions by
\begin{equation}\label{eq:DimlessDef}
  \hat g\smexp{2}=k\smexp{2}\frac{g\smexp{2}_{k}}{f\smexp{2}_{k}},\quad
  \hat y\smexp{2}=y\smexp{2}_{k},\quad
  \xi=\frac{k\smexp{2}}{f\smexp{2}_{k}},\quad
  \kappa=\frac{k\smexp{2}}{M_{\mathrm{Pl}}\smexp{2}(k)}.
\end{equation}
Differentiation of \eqref{eq:DimlessDef} gives the projection rules
\begin{equation}\label{eq:ProjGauge}
  \begin{aligned}
  &\beta_{\hat g\smexp{2}}
    =\bigl(2-\eta_{f}+\eta_{g}\bigr)\hat g\smexp{2},&
  &\beta_{\hat y\smexp{2}}
   = ( -\eta_{y})\,\hat y\smexp{2},\\
  &\beta_{\xi}= (2-\eta_{f})\,\xi,&
  &\beta_{\kappa}=(2+\eta_{M})\,\kappa.
  \end{aligned}
\end{equation}
Hence the task reduces to computing the four anomalous
dimensions \(\{\eta_{f},\,\eta_{g},\,\eta_{y},\,\eta_{M}\}\) from the trace on
the right of the Wetterich equation.


%%%%%%%%%%%%%%%%%%%%%%%%%%%%%%%%%%%%%%%%%%%%%%%%%%%%%%%%%%%%%%%%%%%%%%
\subsection{One-Loop \texorpdfstring{$\beta$}{Beta}-Functions}
\label{app:FRG:OneLoop}
%%%%%%%%%%%%%%%%%%%%%%%%%%%%%%%%%%%%%%%%%%%%%%%%%%%%%%%%%%%%%%%%%%%%%%

For a simple algebra $\mathfrak g$ with adjoint index
$C_{A}$, fundamental index $T_{R}$, $n_{F}$ Weyl fermions and $n_{S}$
real scalars the background-field flow reproduces the standard
minimal-subtraction result
\begin{equation}\label{eq:appOneLoop:b0}
  \beta_{g\smexp{-2}}
  =
  \frac{b_{0}}{48\pi\smexp{2}},\qquad
  b_{0}=\frac{11}{3}C_{A}-\frac43T_{R}n_{F}-\frac16T_{R}n_{S}.
\end{equation}
For the $\ESix$ matter content
\((C_{A},\,T_{R},\,n_{F},\,n_{S})=(12,\,1,\,16,\,32)\), we obtain
$b_{0}=52/3$. The $32$ coset scalars are (pseudo-)Goldstone modes of the breaking
$\ESix\to\SOTenXUOne$; at one loop their two-point function is
protected by the coset symmetry, hence
\begin{equation}\label{eq:appOneLoop:etaf}
  \eta_{f}=-\partial_{t}\log f\smexp{2}_{k}=0.
\end{equation}
$28$ of the $32$ scalars acquire a common mass
\(m\smexp{2}=f\smexp{2}\mu\smexp{2}\) (\autoref{thm:Mass28}). Their one-loop contribution is the functional determinant
\(
  (1/2)\Tr\log\bigl[-D\smexp{2}+m\smexp{2}\bigr].
\)
Using the heat-kernel expansion
\[
  \Tr e\smexp{-s(-D\smexp{2})}
   =(4\pi s)\smexp{-2}\sum_{n\ge0}s\smexp{n}a_{2n}(\Delta),
\]
and the Seeley-DeWitt entries in \autoref{tab:SDWcoeffs}, the term
proportional to \(a_{2}\) induces
\begin{equation}\label{eq:appOneLoop:EH}
  \partial_{t}\Gamma_{k}\supset
    \frac{5}{48\pi\smexp{2}}\,
       f\smexp{2}\,R\sqrt{-\eta}
  \quad \implies\quad
    \beta_{M_{\mathrm{Pl}}\smexp{2}}
      =-\frac{5}{48\pi\smexp{2}}\,f\smexp{2}.
\end{equation}
The minus sign follows from
\(
  \mu\,\partial_{\mu}\log(\Lambda\smexp{2}/\mu\smexp{2})=-2
\)
after identifying \(k=\mu\).
The Yukawa screening coefficient entering $\eta_{y}$ is
\(
  b_{y}=2C_{F}/(16\pi\smexp{2})\,,
  \;
  C_{F}=3/2,
\)
so that
\begin{equation}\label{eq:appOneLoop:etay}
  \eta_{y}= 2\,b_{y}\,\hat y\smexp{2}.
\end{equation}
Collecting \eqref{eq:appOneLoop:b0}--\eqref{eq:appOneLoop:etay}, one finds
\[
  \eta_{f}=0,\quad
  \eta_{g}= \frac{b_{0}}{24\pi\smexp{2}}\hat g\smexp{2},\quad
  \eta_{y}= 2b_{y}\hat y\smexp{2},\quad
  \eta_{M}= -\frac{5}{48\pi\smexp{2}}\kappa.
\]
Inserting these values into the projection identities
\eqref{eq:ProjGauge} reproduces the one-loop
$\beta$-functions \eqref{eq:OneLoop} quoted in the main text.


%%%%%%%%%%%%%%%%%%%%%%%%%%%%%%%%%%%%%%%%%%%%%%%%%%%%%%%%%%%%%%%%%%%%%%
\subsection{Two-Loop \texorpdfstring{$\beta$}{Beta}-Functions}
\label{app:FRG:TwoLoop}
%%%%%%%%%%%%%%%%%%%%%%%%%%%%%%%%%%%%%%%%%%%%%%%%%%%%%%%%%%%%%%%%%%%%%%

Coset-scalar self-interactions and Yukawa diagrams enter at two loops.
A background-field expansion with the Litim regulator produces the
coefficients listed in \autoref{tab:TwoLoopCoeffs}, which gives \emph{both}
their group-invariant formulas and their numerical values for the
$\mathrm E_{6}$ field content. The Mathematica package \texttt{UCFT\_FRG.m} contains the full analytic derivation and writes these numbers to \texttt{twoLoopCoeffs.dat}.

\begin{table}[h]
\renewcommand{\arraystretch}{1.5}
\centering
\caption{Two-loop $\beta$-function coefficients for the $\mathrm E_{6}$ field content $(C_{A}=12,\;C_{F}=9/4,\;T_{R}=1,\;n_{F}=16,\;n_{S}=32)$.}
\label{tab:TwoLoopCoeffs}
\renewcommand{\arraystretch}{1.35}
\begin{tabular}{lcc}
\hline
Coefficient & Symbolic expression & $\mathrm E_{6}$ value \\ \hline
$\displaystyle A$   &
\rule{0pt}{1.8em} $\dfrac{T_{R}}{192\,\pi\smexp{4}}$ &
$5.35\times10\smexp{-5}$ \\[9pt]

$\displaystyle b_{0}$ &
$\dfrac{\;\tfrac{11}{3}C_{A}\;-\;\tfrac{4}{3}T_{R}n_{\sm F}\;-\;\tfrac16 T_{R}n_{\sm S}}{16\,\pi\smexp{2}}$ &
$1.10\times10\smexp{-1}$ \\[9pt]

$\displaystyle b_{1}$ &
$\dfrac{34\,C_{A}\smexp{2}\;-\;4\,C_{A}T_{R}n_{\sm F}\;-\;\tfrac53 C_{A}T_{R}n_{\sm S}\;-\;4\,C_{F}T_{R}n_{\sm F}}
       {(16\,\pi\smexp{2})\smexp{2}}$ &
$1.34\times10\smexp{-1}$ \\[9pt]

$\displaystyle c_{1}$ &
$-\dfrac{3\,C_{F}}{16\,\pi\smexp{2}}$ &
$-4.27\times10\smexp{-2}$ \\[9pt]

$\displaystyle c_{2}$ &
$\dfrac{2\,C_{F}\;-\;\tfrac32\,C_{A}}{16\,\pi\smexp{2}}$ &
$-8.55\times10\smexp{-2}$ \\[9pt]

$\displaystyle b_{y}$ &
$\dfrac{2\,C_{F}}{(16\,\pi\smexp{2})}$ &
$2.85\times10\smexp{-2}$ \\[9pt] \hline
\end{tabular}
\end{table}

%%%%%%%%%%%%%%%%%%%%%%%%%%%%%%%%%%%%%%%%%%%%%%%%%%%%%%%%%%%%%%%%%%%%%%
\subsection{Stability Matrix at the Fixed Point}
\label{app:FRG:Stability}
%%%%%%%%%%%%%%%%%%%%%%%%%%%%%%%%%%%%%%%%%%%%%%%%%%%%%%%%%%%%%%%%%%%%%%

Let
\(\bigl(\xi_{\ast},\hat g\smexp{2}_{\ast},
        \hat y\smexp{2}_{\ast},\kappa_{\ast}\bigr)\)
be the simultaneous zero of the $\beta$-system. Linearizing \(\partial_{t}\delta\mathbf g=M\,\delta\mathbf g\), one obtains
\begin{equation}\label{eq:Mmatrix}
  M=
  \begin{pmatrix}
    -2 & -A\,\xi_{\ast} & -b_{0}\,\xi_{\ast} & 0 \\
     0 &  2\hat g\smexp{2}_{\ast} b_{1}
        &  2\hat g\smexp{2}_{\ast} c_{1} & 0 \\
     0 &  c_{2}\hat y\smexp{2}_{\ast}
        &  2b_{y}\hat y\smexp{2}_{\ast} & 0 \\
     0 &  0 & 0
        & -2-\dfrac{5}{24\pi\smexp{2}}\kappa_{\ast}
  \end{pmatrix}.
\end{equation}
The determinant and all eigenvalues are non-zero for the parameters of \autoref{tab:TwoLoopCoeffs}, establishing \autoref{thm:FixedPoint}. Evaluating \eqref{eq:Mmatrix} at the fixed point gives \(\sigma(M)=\{-2.0,-1.3,-0.9,-0.2\}\).

%%%%%%%%%%%%%%%%%%%%%%%%%%%%%%%%%%%%%%%%%%%%%%%%%%%%%%%%%%%%%%%%%%%%%%
\subsection{Finite Coset-Scalar Four-Point Amplitude}
\label{app:FRG:4pt}
%%%%%%%%%%%%%%%%%%%%%%%%%%%%%%%%%%%%%%%%%%%%%%%%%%%%%%%%%%%%%%%%%%%%%%

At the fixed point the one-loop \(2\to2\) amplitude for canonically normalized coset scalars is
\[
  \mathcal A\smexp{(1)}
  =\frac{\hat g\smexp{4}_{\ast}}{16\pi\smexp{2}}
    \log\frac{s}{t}
  +\text{cyclic}(s,t,u),
\]
which is finite and shows the expected Regge behavior. See \cite{Percacci2017} for method details.

%======================================================================%
\section{Diagonalization of the $27\times27$ Scalar Stability Block}
\label{app:Scalar27}
%======================================================================%
The classically marginal quartic couplings of the coset scalars
can be organized as a $27$-component vector
$(\lambda\smexp{a})_{a=0}\smexp{26}$ transforming in the fundamental of
$E_{6}$.  Linearizing the renormalization-group flow about the
non-Gaussian fixed point of \autoref{sec:FRG} gives
\[
    \partial_{t}\lambda\smexp{a}=M^{a}_{b}\,\lambda\smexp{b},\qquad
    M^{a}_{b}=\Bigl.\frac{\partial\beta\smexp{a}}{\partial\lambda\smexp{b}}
    \Bigr|_{\ast}.
\]

Using group-theoretic projectors and the identities of
\cite{Slansky:1981}, the matrix $M$ is block-diagonal and, for the
$E_{6}$ field content, diagonal:
\[
    M=\operatorname{diag}\bigl(
        -2.00,\,-1.89,\,-1.57,\,-1.18,\,-0.92,\,-0.77,\,-0.44
        \ (\times21)\bigr),
\]
all eigenvalues being negative.  Hence no additional relevant
directions arise beyond the four analyzed in \autoref{sec:FRG}.

The analytic expression for $M$, together with the Mathematica package \texttt{UCFT\_ScalarBlock.m}, is provided in the ancillary files.


\end{appendices}

%%===========================================================================================%%
%% If you are submitting to one of the Nature Portfolio journals, using the eJP submission   %%
%% system, please include the references within the manuscript file itself. You may do this  %%
%% by copying the reference list from your .bbl file, paste it into the main manuscript .tex %%
%% file, and delete the associated \verb+\bibliography+ commands.                            %%
%%===========================================================================================%%

\bibliography{sn-bibliography}

\end{document}
